\section{Perfumes}

\subsection{Presentación del caso}
Una empresa cosmética elabora tres tipos de perfumes: perfume floral, perfume amaderado y perfume cítrico.
Las contribuciones unitarias al beneficio son, respectivamente, 4, 6 y 5 unidades monetarias por m$^3$ producido.

La producción se organiza en tres procesos principales: maceración de esencias, mezclado y envasado, que requieren la dedicación de 1, 2 y 4 trabajadores respectivamente.

Cada tipo de perfume emplea los siguientes tiempos de proceso (en minutos por m$^3$):

\begin{center}
\begin{tabular}{lccc}
\toprule
 & Maceración & Mezclado & Envasado \\
\midrule
Perfume floral & 4 & 6 & 8 \\
Perfume amaderado & 4 & 10 & 15 \\
Perfume cítrico & 4 & 12 & 5 \\
\bottomrule
\end{tabular}
\end{center}

Los tiempos máximos disponibles al día son 240 minutos de maceración, 480 minutos de mezclado y 960 minutos de envasado.

Se pide plantear un modelo de programación lineal que maximice el beneficio.

\subsection{Formulación explícita}

\paragraph{Variables}\mbox{}\\
\begin{tabular}{lcl}
$x_f$ & : & cantidad (m$^3$) de perfume floral producido \\
$x_a$ & : & cantidad (m$^3$) de perfume amaderado producido \\
$x_c$ & : & cantidad (m$^3$) de perfume cítrico producido \\
\end{tabular}

\paragraph{Restricciones}\mbox{}\\

Tiempo máximo de maceración:
\begin{equation}
4x_f + 4x_a + 4x_c \leq 240
\end{equation}

Tiempo máximo de mezclado:
\begin{equation}
6x_f + 10x_a + 12x_c \leq 480
\end{equation}

Tiempo máximo de envasado:
\begin{equation}
8x_f + 15x_a + 5x_c \leq 960
\end{equation}

\paragraph{Función objetivo}\mbox{}\\
\begin{equation}
\mbox{max.}\; z = 4x_f + 6x_a + 5x_c
\end{equation}

\paragraph{Modelo completo}\mbox{}\\
\begin{align}
\mbox{max.}\quad & 4x_f + 6x_a + 5x_c \notag \\
\mbox{s.a.}\quad
& 4x_f + 4x_a + 4x_c \leq 240 \notag \\
& 6x_f + 10x_a + 12x_c \leq 480 \notag \\
& 8x_f + 15x_a + 5x_c \leq 960 \notag \\
& x_f, x_a, x_c \geq 0
\end{align}

\subsection{Formulación generalizada}

\paragraph{Conjuntos}\mbox{}\\
\begin{tabular}{lcl}
$\mathcal{F}$ & : & tipos de perfumes \\
$\mathcal{P}$ & : & procesos de producción (maceración, mezclado, envasado) \\
\end{tabular}

\paragraph{Parámetros}\mbox{}\\
\begin{tabular}{lcl}
$B_f$ & : & beneficio unitario por m$^3$ de perfume $f \in \mathcal{F}$ \\
$T_{pf}$ & : & tiempo requerido (min) del proceso $p \in \mathcal{P}$ por cada m$^3$ de perfume $f \in \mathcal{F}$ \\
$D_p$ & : & tiempo disponible (min) para el proceso $p \in \mathcal{P}$ \\
\end{tabular}

\paragraph{Variables}\mbox{}\\
\begin{tabular}{lcl}
$x_f$ & : & cantidad de perfume $f \in \mathcal{F}$ a producir (m$^3$) \\
\end{tabular}

\paragraph{Restricciones}\mbox{}\\
Para cada proceso $p \in \mathcal{P}$:
\begin{equation}
\sum_{f \in \mathcal{F}} T_{pf}\,x_f \leq D_p
\end{equation}

\paragraph{Función objetivo}\mbox{}\\
\begin{equation}
\mbox{max.}\; z = \sum_{f \in \mathcal{F}} B_f\,x_f
\end{equation}

\paragraph{Modelo completo}\mbox{}\\
\begin{align}
\mbox{max.}\quad & \sum_{f \in \mathcal{F}} B_f\,x_f \notag \\
\mbox{s.a.}\quad
& \sum_{f \in \mathcal{F}} T_{pf}\,x_f \leq D_p, \quad \forall p \in \mathcal{P} \notag \\
& x_f \geq 0, \quad \forall f \in \mathcal{F}
\end{align}
